\begingroup
\makeatletter\tagsleft@true\let\veqno\@@leqno\makeatother
\ArithmeticHead{349}{R}{{\scriptsize ZUR THEORIE DER REIHE $F(\alpha,\beta,\gamma,x)$.\quad III.}}
\NumberedItem{7}
\noindent{\small[S. 45]}\par
\indent Wir kommen auf die Gleichung (30) zurück, und bemerken, dass wenn allgemein
\[F(\alpha,\beta,\gamma,x)=P\ \text{und}\ F(\alpha,\beta,\alpha+\beta+1-\gamma,1-x)=Q\]
\noindent gesetzt wird, die beiden Gleichungen Statt haben
\[\begin{aligned}
\alpha\beta P-(\gamma-(\alpha+\beta+1)x)\frac{dP}{dx}-(x-xx)\frac{d^2P}{dx^2}&=0,\\[.6em]
\alpha\beta Q-(\gamma-(\alpha+\beta+1)x)\frac{dQ}{dx}-(x-xx)\frac{d^2Q}{dx^2}&=0.
\end{aligned}\]
\noindent Hieraus folgt \textls[120]{erstlich}, dass das vollständige Integral der Gleichung
\[\alpha\beta y-(\gamma-(\alpha+\beta+1)x)\frac{dy}{dx}-(x-xx)\frac{d^2y}{dx^2}=0\]
\noindent sey
\[y=\mathfrak{A}F(\alpha,\beta,\gamma,x)+\mathfrak{B}F(\alpha,\beta,\alpha+\beta+1-\gamma,1-x),\]
\noindent wo $\mathfrak{A}$ und $\mathfrak{B}$ Constanten sind.
\par\smallskip\noindent{\small[S. 46]}\par
\indent\textls[120]{Zweitens}, dass
\[-\frac{\gamma-(\alpha+\beta+1)x}{x-xx}dx=d\log(QdP-PdQ),\]
\noindent also
\begin{equation}
\frac{QdP-PdQ}{dx}=\frac{A}{x^\gamma(1-x)^{\alpha+\beta+1-\gamma}}\mathpunct{\raisebox{.35ex}{.}}\tag*{[31a]}
\end{equation}
\clearpage\SourceLeaf{357}\ArithmeticHead{350}{L}{ANALYSIS.\quad NACHLASS.}
\par\indent Noch eine wichtige Verwandlung vollführt man auf folgende Art:
\par\indent Man setze
\[x=\frac{t}{t-1},\]
\noindent also
\[\frac{dx}{dt}=-\frac{1}{(t-1)^2},\qquad\frac{d^2x}{dt^2}=\frac{2}{(t-1)^3}\mathpunct{\raisebox{.35ex}{.}}\]
\noindent Hiedurch verwandelt sich die Gleichung für $F(\alpha,\beta,\gamma,x)=P$ in
\[\frac{\alpha\beta}{t-1}P-(\gamma+(\alpha+\beta-1-\gamma)t)\frac{dP}{dt}-(t-tt)\frac{d^2P}{dt^2}=0.\]
\noindent Macht man also
\[P=(t-1)^\mu P',\]
\noindent so wird
\[\begin{aligned}
&\frac{\alpha\beta-\gamma\mu+\mu t(\gamma-\alpha-\beta+\mu)}{t-1}P'\\
&\qquad-(\gamma+(\alpha+\beta-1-\gamma-2\mu)t)\frac{dP'}{dt}-(t-tt)\frac{d^2P'}{dt^2}=0.
\end{aligned}\]
\noindent Um den Bruch wegzuschaffen, setzt man
\[\alpha\beta-\alpha\mu-\beta\mu+\mu\mu=0,\]
\noindent also entweder $\mu=\alpha$ oder $\mu=\beta$. Im ersten Falle wird
\[(-\alpha\beta+\alpha\gamma)P'-(\gamma-(\gamma+1+\alpha-\beta)t)\frac{dP'}{dt}-(t-tt)\frac{d^2P'}{dt^2}=0,\]
\noindent woraus
\[P'=\mathfrak{A}'F(\alpha,\gamma-\beta,\gamma,t)+\mathfrak{B}'F(\alpha,\gamma-\beta,\alpha+1-\beta,1-t).\]
\noindent Hieraus folgert man leicht
\[\begin{aligned}
F(\alpha,\beta,\gamma,x)&=(1-x)^{-\alpha}F\left(\alpha,\gamma-\beta,\gamma,-\frac{x}{1-x}\right),\\[.7em]
&=(1-x)^{-\beta}F\left(\beta,\gamma-\alpha,\gamma,-\frac{x}{1-x}\right).
\end{aligned}\]
\noindent Mit Hülfe dieser Verwandlung kann man die Werthe unserer Functionen für negative $x$ auf die Werthe für positive reduciren.[*)]
\[\begin{aligned}
F\left(\alpha,\beta,\gamma,\frac1x\right)&=Ax^\alpha F(\alpha,\alpha+1-\gamma,\alpha+1-\beta,x)\\[.7em]
&\qquad+Bx^\beta F(\beta,\beta+1-\gamma,\beta+1-\alpha,x),
\end{aligned}\]
\par\medskip\noindent\rule{.25\linewidth}{.4pt}\par\smallskip
{\small\noindent [*) Die folgenden Formeln ohne Text sind ersichtlich später hinzugefügt.]}
\clearpage\SourceLeaf{358}\ArithmeticHead{351}{R}{{\scriptsize ZUR THEORIE DER REIHE $F(\alpha,\beta,\gamma,x)$.\quad III.}}
\[\begin{aligned}
A&=(\cos\alpha\pi+i\sin\alpha\pi)\frac{\Pi(\gamma-1)\Pi(\beta-1-\alpha)}{\Pi(\beta-1)\Pi(\gamma-1-\alpha)}\\[1em]
B&=(\cos\beta\pi+i\sin\beta\pi)\frac{\Pi(\gamma-1)\Pi(\alpha-1-\beta)}{\Pi(\alpha-1)\Pi(\gamma-1-\beta)}\mathpunct{\raisebox{.35ex}{.}}
\end{aligned}\]
\clearpage
\endgroup
